\begin{center}
	\Large
	Uden hjælpemidler	
\end{center}
%
%
%
%
%
\begin{opgavetekst}{Opgave 1}
\end{opgavetekst}
%
%
\begin{delopgave}{1}
	Udregn følgende.
	\begin{align*}
		&1) \ 41 \fmod{9}   &&2) \ 77 \fmod{11}   &&3) \ 12 \fmod{2}   \\
		&4) \ -57 \fmod{15}   &&5) \ 22 \fmod{19}   &&6) \ 59^{6} \fmod{3}   \\
	\end{align*}
\end{delopgave}
%
%
%
%
%
\begin{opgavetekst}{Opgave 2}
\end{opgavetekst}
%
%
\begin{delopgave}{1}
	Bestem en multiplikationstabel for $\mathbb{Z}/11\mathbb{Z}$. 
\end{delopgave}
%
%
\begin{delopgave}{2}
	Afgør, hvilke elementer der er multiplikative inverse i $\mathbb{Z}/11\mathbb{Z}$. 
\end{delopgave}
%
%
%
%
%
\begin{opgavetekst}{Opgave 3}
\end{opgavetekst}
\begin{delopgave}{1}
	Skriv følgende tal i base 10.
	\begin{align*}
		&1) \ 101101_2   &&2) \ 221_3   &&3) \ 414_5    \\
		&4) \ 65_7   &&5) \ 333_4    &&6) \ 55_6   \\
	\end{align*}
\end{delopgave}
\begin{delopgave}{2}
	Omskriv følgende tal fra base 10 til den angivne base.
	\begin{align*}
		&1) \ 100\textnormal{ i base 2}   &&2) \ 37\textnormal{ i base 5}   &&3) \ 69\textnormal{ i base 4}    \\
		&4) \ 99\textnormal{ i base 9}   &&5) \ 64\textnormal{ i base 8}    &&6) \ 242\textnormal{ i base 3}   \\
	\end{align*}	 
\end{delopgave}
%
%
%
%
%
\begin{opgavetekst}{Opgave 4}
\end{opgavetekst}
\begin{delopgave}{1}
	Udregn følgende
	\begin{align*}
		&1) \ 331_4 + 133_4 &&2) \ 51_5\cdot 103_5
	\end{align*}
\end{delopgave}
%
%
%
%
%
\newpage
%
%
%
%
%
\begin{opgavetekst}{Opgave 5}
\end{opgavetekst}
\begin{delopgave}{1}
	Udregn største fælles divisor for følgende tal uden at bruge Euklids algoritme.
	\begin{align*}
		&1) \ \textnormal{gcd}(12,18)   &&2) \ \textnormal{gcd}(25,50)   &&3) \ \textnormal{gcd}(11,17)    \\
		&4) \ \textnormal{gcd}(26,38)   &&5) \ \textnormal{gcd}(100,101)    &&6) \ \textnormal{gcd}(77,70)   \\
	\end{align*}
\end{delopgave}
%
%
\begin{delopgave}{2}
	Udregn største fælles divisor for følgende tal ved at bruge Euklids algoritme
	\begin{align*}
		&1) \ \textnormal{gcd}(130,56)   &&2) \ \textnormal{gcd}(201,21)   &&3) \ \textnormal{gcd}(18,157)    \\
	\end{align*}
\end{delopgave}
%
%
\begin{delopgave}{3}
	\begin{enumerate}[label=\roman*)]
		\item Bestem den multiplikative inverse til $13$ i $\mathbb{Z}/57\mathbb{Z}$.
		\item Bestem den multiplikative inverse til $199$ i $\mathbb{Z}/25\mathbb{Z}$.
	\end{enumerate}
\end{delopgave}
%
%
%
%
%
\begin{opgavetekst}{Opgave 6}
\end{opgavetekst}
%
%
\begin{delopgave}{1}
	Udregn følgende
	\begin{align*}
		&1) \ \varphi(11)   &&2) \ \varphi(15)   &&3) \ \varphi(16)
	\end{align*}	 
\end{delopgave}
%
%
%
%
%
\begin{opgavetekst}{Opgave 7}
	Brug Fermats lille sætning til at udregne følgende
	\begin{align*}
		&1) \ 25^{62} \fmod{7}   &&2) \ 4^{1003} \fmod{11}   &&3) \ 11^{1025} \fmod{3} \\
		&4) \ 22^{31} \fmod{31}   &&5) \ 22^{48} \fmod{5}   &&6) \ 77319412^{58} \fmod{59}
	\end{align*}	 
\end{opgavetekst}